\chapter{OBJECTIVES AND SYSTEM REQUIREMENTS}

\section{Requirement Basis}
System requirements were formulated based on the core engineering objectives of the Real-Time Wellness Mirror, integrating FreeRTOS determinism, secure telemetry ingestion, rigorous schema validation, and sensor-grounded conversational AI. The specifications balance functional richness with edge execution constraints, privacy, and non-clinical safety boundaries.

\section{Functional Requirements}
The functional requirements define the operational capabilities that the system must provide across its sensing, task scheduling, bridge validation, and interaction subsystems. Table~\ref{tab:functional-requirements} records the functional requirements, the adopted implementation approaches, and their verified demonstration status.

\begin{table}[H]
\centering
\scriptsize
\renewcommand{\arraystretch}{0.98}
\caption{Functional requirements and demonstration verification status.}
\label{tab:functional-requirements}
\begin{tabularx}{\textwidth}{p{8mm}p{28mm}p{42mm}p{34mm}Y}
\toprule
\textbf{ID} & \textbf{Requirement} & \textbf{Implementation Approach} & \textbf{Acceptance Criteria} & \textbf{Status}\\
\midrule
FR1 & Task Scheduling & FreeRTOS kernel with tasks, queue (8 slots), and mutex. & Preemptive execution with observed lateness and queue HWM. & Implemented and tested\\
FR2 & Multi-Modal Acquisition & Android accessibility relay for wearable UI; RTOS simulator. & Ingest HR, $\mathrm{SpO_2}$, steps from relay; synthetic RTOS stream. & Partially verified\\
FR3 & Telemetry Validation & Ingress validator checking 8 keys, types, physical ranges. & Reject incomplete, simulated, out-of-range, or malformed packets. & Implemented and tested\\
FR4 & Freshness Tracking & Web layer tracking arrival age and snapshot fingerprint. & Visual indicator transition: Live, Unchanged, Stale, Offline ($>300$\,s). & Implemented and tested\\
FR5 & Grounded AI Guidance & Local Ollama \code{gemma2:2b} with rule-based consistency check. & Response bounded to 70 words; explicit \texttt{PASS}/\texttt{FAIL} verdict. & Partially verified\\
FR6 & Natural Voice Input & Local faster-whisper speech-to-text (\code{tiny.en}, \code{int8}). & Low-latency local transcription and audible speech response. & Partially verified\\
FR7 & Standalone Fallback & Standalone RTOS synthetic mode with console controls. & Continuous operation when external wearable transport is absent. & Implemented and tested\\
FR8 & Physical ESP32 Node & ESP32 pinout with MPU6050, DHT22, HC-SR04 sensors. & Real-time physical sensor acquisition on embedded hardware. & Planned / Simulation\\
\bottomrule
\end{tabularx}
\end{table}

\section{Non-Functional Requirements}
Non-functional requirements define the operational qualities, architectural constraints, and safety boundaries necessary for a reliable, responsive, and privacy-preserving smart mirror. Table~\ref{tab:nonfunctional-requirements} lists these specifications alongside their control mechanisms and verification status.

\begin{table}[H]
\centering
\scriptsize
\renewcommand{\arraystretch}{0.98}
\caption{Non-functional requirements and engineering specifications.}
\label{tab:nonfunctional-requirements}
\begin{tabularx}{\textwidth}{p{9mm}p{30mm}p{56mm}Y}
\toprule
\textbf{ID} & \textbf{Requirement} & \textbf{Design Specification and Control Mechanism} & \textbf{Status}\\
\midrule
NFR1 & Real-Time Observability & Bounded release monitoring in FreeRTOS with queue HWM, dropped-packet, and missed-deadline tracking. & Partially verified (host simulation)\\
NFR2 & Telemetry Freshness & Explicit arrival age tracking with automated visual blanking of metrics when data exceeds the 300\,s offline threshold. & Implemented and tested\\
NFR3 & Host Privacy \& Locality & Ingest telemetry and execute speech-to-text (faster-whisper) and language inference (local Ollama gemma2:2b) entirely on the host workstation without external cloud routing; relay transport uses authenticated USB port forwarding or private Wi-Fi. & Implemented and tested\\
NFR4 & Data Provenance Isolation & Enforce strict separation between genuine wearable data and synthetic control streams to prevent data conflation. & Implemented and tested\\
NFR5 & Interface Robustness & Handle transport disconnects, malformed JSON, and unexpected key sequences gracefully without browser stall. & Implemented and tested\\
NFR6 & Non-Clinical Safety & Restrict assistant advice to non-diagnostic wellness feedback with prominent clinical disclaimers. & Implemented and tested\\
\bottomrule
\end{tabularx}
\end{table}

\section{Shared Telemetry Contract}
The wellness mirror dashboard and assistant ingest telemetry exclusively through the locked eight-field contract shown in Table~\ref{tab:telemetry-contract}. The contract enforces semantic constraints, ensuring that derived heuristics and compatibility values are clearly distinguished from direct physiological measurements before reaching the presentation surface.

\begin{table}[H]
\centering
\scriptsize
\renewcommand{\arraystretch}{0.98}
\caption{Eight-field telemetry contract enforced by the shared receiver.}
\label{tab:telemetry-contract}
\begin{tabularx}{\textwidth}{p{28mm}p{22mm}p{44mm}Y}
\toprule
\textbf{Field} & \textbf{Data Type} & \textbf{Validation \& Acceptance Range} & \textbf{Operational Interpretation}\\
\midrule
\code{timestamp} & ISO-8601 string & Required; valid date-time. & Telemetry generation or capture timestamp.\\
\code{heart\_rate} & Integer BPM & Required; $30 \le \mathrm{HR} \le 220$. & Directly observed or simulated heart rate.\\
\code{spo2} & Integer percent & Required; $70 \le \mathrm{SpO_2} \le 100$. & Peripheral capillary oxygen saturation.\\
\code{activity} & Enum string & \code{rest}, \code{walk}, \code{run}, \code{active}, \code{unknown}. & Compatibility default is \code{unknown} if unmeasured.\\
\code{motion\_intensity} & Float & Required; $0.0 \le m \le 1.0$. & Heuristic proxy; $0.0$ does not imply total immobility.\\
\code{exertion} & Enum string & \code{rest}, \code{light}, \code{moderate}, \code{high}. & Derived heuristic estimate from HR and step rate.\\
\code{steps} & Integer & Required; non-negative integer $\ge 0$. & Cumulative daily step count.\\
\code{ppg\_quality} & Enum string & \code{good} or \code{poor}. & Upstream completion proxy; not a raw PPG algorithm.\\
\bottomrule
\end{tabularx}
\end{table}

\section{Safety Boundaries and Acceptance Criteria}
The accepted telemetry context represents transport data, not a clinical guarantee. A grounding \texttt{PASS} indicates that the generated assistant text contains required contextual mentions consistent with the active snapshot; a \texttt{FAIL} indicates that one or more checks failed. Neither verdict implies medical correctness. Under the current implementation, speech synthesis continues following a failed check, serving as an informative indicator rather than a fail-closed interlock. Emergency response, diagnosis, and medical treatment are explicitly excluded from system scope.

\section{Repeatable Demonstration Sequence}
The formal demonstration sequence is structured into six repeatable stages:
\begin{enumerate}[noitemsep,topsep=1pt]
  \item Initialize the local mirror application on the host laptop (\code{http://127.0.0.1:8777}) and start the ingestion receiver on port 8000;
  \item Establish relay transport via USB reverse port forwarding (\code{adb reverse tcp:8000 tcp:8000}) or private Wi-Fi IP, and start the Android streamer;
  \item Inspect the mirror display to confirm active link state and source provenance labels;
  \item Transmit a verified contract payload and observe real-time telemetry card updates;
  \item Invoke hands-free or text assistance to generate a bounded wellness response and review the grounding verdict; and
  \item Simulate transport disconnection to verify transitions to \texttt{STALE} and subsequent metric blanking in \texttt{OFFLINE}.
\end{enumerate}